\documentclass[10pt]{article}

\usepackage[utf8]{inputenc}

\usepackage[T1]{fontenc}

\usepackage{amsmath}

\usepackage{amsfonts}

\usepackage{amssymb}

\usepackage[version=4]{mhchem}

\usepackage{stmaryrd}



\begin{document}

$\in L_{2}[0, \pi]$ имеет место так наз. $\mathrm{p}$ а в е в с т в о $П$ а $\mathrm{p}-$ с е в а л я



\[

\int_{0}^{\pi}|f(x)|^{2} d x=\sum_{n=0}^{\infty}\left|a_{n}\right|^{2}

\]



где



\[

a_{n}=\int_{0}^{\pi} f(x) v_{n}(x) d x

\]



и справедлива формула разложения по собственным функциям



\[

f(x)=\sum_{n=0}^{\infty} a_{n} v_{n}(x),

\]



где ряд сходится в метрике пространства $L_{2}[0, \pi]$. Теоремы полноты и разложения для регулярной Ш.Л. з. впервые доказаны В. А. Стекловым [14].



Если функция $f(x)$ имеет вторую непрерывную производную и удовлетворяет граничным условиям (3), то справедливы следующие утверждения (см. [15]):



a) ряд (4) сходится абсолютно в равномерно на сегменте $[0, \pi]$ к функции $f(x)$;



б) один раз продифферевцированный ряд (4) сходится абсолютно и равномерно на сегменте $[0, \pi]$ к $f^{\prime}(x)$;



в) в каждой точке, в к-рой $f^{\prime \prime}(x)$ удовлетворяет какому-либо локальному условию разложения в ряд Фурье (вапр., имеет огравиченную вариацию), дважды продифференцированный ряд (4) сходится к $f^{\prime \prime}(x)$.



Для любой функции $f(x) \in L_{1}[0, \pi]$ ряд (4) является равномерно равносходящимся с рядом Фурье фуункции $f(x)$ по $\cos n x$, т. е.



\[

\lim _{N \rightarrow \infty} \sup _{0 \leqslant x \leqslant \pi}\left|V_{N, f}(x)-c_{N, f}(x)\right|=0,

\]



где



\[

\begin{gathered}

V_{N, f}(x)=\int_{0}^{\pi} f(t) \sum_{n=0}^{N} v_{n}(x) v_{n}(t) d t, \\

c_{N, f}(x)=\int_{0}^{\pi} f(t)\left\{\frac{1}{\pi}+\frac{2}{\pi} \sum_{n=1}^{N} \cos n x \cos n t\right\} d t .

\end{gathered}

\]



Это утверждение означает, что разложение функции $f(x)$ по собственным функциям граничной задачи (2), (3) сходится при тех же условиях, что и разложение $f(x)$ в ряд Фурье по косинусам (см. [1], [4]).



\begin{enumerate}

  \setcounter{enumi}{1}

  \item Рассматривается дифференциальное уравнение (2) на полуоси $0 \leqslant x<\infty$ с граничным условием в нуле:

\end{enumerate}



\[

y^{\prime}(0)-h y(0)=0 .

\]



Функция $q(x)$ предполагается действительной и суммируемой в каждом конечном подинтервале интервала $[0, \infty)$, а число $h$ действительным.



Пусть $\varphi(x, \lambda)$ - решение уравнения (2) с начальными условиями $y(0)=1, y^{\prime}(0)=h$ (так что $\varphi(x, \lambda)$ удовлетворяет и граничному условию (5)). Пусть $f(x)-$ любая функция из $L_{2}(0, \infty)$ и $\Phi_{f, b}(x)=\int_{0}^{b} f(x) \varphi(x, \lambda) d x$, где $b$ - произвольное конечное положительное число. Для каждой функции $q(x)$ и каждого числа $h$ существует, по крайней мере, одна, не зависящая от $f(x)$, неубывающая функция $\rho(\lambda),-\infty<\lambda<\infty$, обладающая следующими свойствами:



a) существует функция $\Phi_{f}(\lambda)$, являющаяся пределом $\Phi_{f, b}(\lambda)$ при $b \rightarrow \infty$ в метрике $L_{2, \rho}(-\infty, \infty)$ (пространства $\rho$-измеримых функций $\psi(\lambda)$, для к-рых



\[

\left.\|\psi\|=\int_{-\infty}^{\infty}|\psi(\lambda)|^{2} d \rho(\lambda)<\infty\right),

\]



T. e.



\[

\lim _{b \rightarrow \infty} \int_{-\infty}^{\infty}\left|\Phi_{f}(\lambda)-\Phi_{f, b}(\lambda)\right|^{2} d \rho(\lambda)=0 ;

\]



б) имеет место равенство Парсеваля



\[

\int_{0}^{\infty}|f(x)|^{2} d x=\int_{-\infty}^{\infty}\left|\Phi_{f}(\lambda)\right|^{2} d \rho(\lambda)

\]



Функция $\rho(\lambda)$ наз. с п е к т р а л b в 0 й ци е й (или с п е к т рально й пл от т ос т ь ю) граничной задачи (2), (5) (см. [9]-[11]).



Для спектральной функции $\rho(\lambda)$ задачи $(2),(5)$ справедлива асимптотич. формула (см. [16]) (в уточненном виде см. [17]):



\[

\begin{gathered}

\lim _{\lambda \rightarrow-\infty} e^{\sqrt{\lambda} x}(\rho(\lambda)-\rho(-\infty))=0,0 \leqslant x<\infty, \\

\lim _{\lambda \rightarrow \infty}\left(\rho(\lambda)-\rho(-\infty)-\frac{2}{\pi} \sqrt{\bar{\lambda}}+h\right)=0 .

\end{gathered}

\]



Справедлива следующая теорема равносходимости: для произвольной функции $f(x) \in L_{2}(0, \infty)$ пусть



\[

\begin{gathered}

\Phi_{f}(\lambda)=\int_{0}^{\infty} f(x) \varphi(x, \lambda) d x, \\

C_{f}(\lambda)=\int_{0}^{\infty} f(x) \cos \sqrt{\lambda} x d x

\end{gathered}

\]



(интегралы сходятся в мөтриках пространств $L_{2, \rho}(-\infty, \infty)$ и $L_{2,} \sqrt{\lambda}^{-}(0, \infty)$ соответственно); тогда при каждом фиксированном $b<\infty$ сходится интеграл



\[

\int_{-\infty}^{0} \Phi_{f}(\lambda) \varphi(x, \lambda) d \rho(\lambda)

\]



абсолютно и равномерно относительно $x \in[0, b]$ и



\[

\begin{gathered}

\lim _{N \rightarrow \infty} \sup _{0 \leqslant x<b} \mid \int_{-\infty}^{N} \Phi_{f}(\lambda) \varphi(x, \lambda) d \rho(\lambda)- \\

-\frac{2}{\pi} \int_{0}^{N} C_{f}(\lambda) \cos \sqrt{\lambda} x d \sqrt{\lambda} \mid=0 .

\end{gathered}

\]



Пусть задача (2), (5) имеет дискретный спектр, т. е. ее спектр состоит из счетного тисла собственных значевий $\lambda_{1}, \lambda_{2}, \ldots,<\lambda_{n},<\ldots$ с единственной предельной точкой в бесконечности. При определенных условиях на функцию $q(x)$ для функции $N(\lambda)=\Sigma_{\lambda_{n}<\lambda} 1$, т. е. числа собственных значений, меньших $\lambda$, справедлива асимптотич. формула:



\[

N(\lambda) \sim \frac{1}{2 \pi} \int_{q(x)<\lambda}(\lambda-q(x))^{1 / 2} d x .

\]



Наряду с решевием $\varphi(x, \lambda)$ вводится второе решение $\theta(x, \lambda)$ уравнения (2), удовлетворяющее условиям $\theta(0, \lambda)=0, \theta^{\prime}(0, \lambda)=1$, так что $\varphi(x, \lambda)$ и $\theta(x, \lambda)$ образуют фундаментальную систему решевий уравнения (2). При фиксированных числах $\lambda(\operatorname{Im} \lambda \neq 0)$ и $b>0$ рассматривается дробно-линейная функция



\[

w_{\lambda, b}=w_{\lambda, b}(t)=\frac{-\theta^{\prime}(b, \lambda)-t \theta(b, \lambda)}{\varphi^{\prime}(b, \lambda)+t \varphi(b, \lambda)} .

\]



Когда независимая переменная $t$ пробегает действительную ось, точка $w \lambda, b$ описывает нек-рую окружность, ограничивающую круг $K \lambda, b$. Он всегда лежит в той же полуплоскости (нижней или верхней), что и $\lambda$. С увеличением $b$ круг $K \lambda, b$ сжимается, т. е. при $b<b^{\prime}$ круг $K \lambda, b^{\prime}$ лежит деликом внутри круга $K \lambda, b$. Существует (при $b \rightarrow \infty$ ) предельный круг или точка $K_{\lambda, \infty} ;$ при этом если



\[

\int_{0}^{\infty}|\varphi(x, \lambda)|^{2} d x<\infty,

\]



то $K_{\lambda, \infty}$ будет кругом, и точкой - в противном случае (см. [10]). Если условие (6) выполняется для одного какого-либо недействительного значения $\lambda$, то оно выполняется для всех значений $\lambda$. В случае предельного круга для всех значений $\lambda$ все решения уравнения (2) принадлежат пространству $L_{2}(0, \infty)$, а в случае предельной точки для каждого недействительного значения $\lambda$ это уравнение имеет решение вида $\theta(x, \lambda)+m(\lambda) \varphi(x, \lambda)$, принадлежащее $L_{2}(0, \infty)$, где $m(\lambda)$ - предельная точка $\left(m(\lambda)=K_{\lambda, \infty}\right)$.



Если $q(x) \geqslant-c x^{2}$, где $c$ - нек-рая положительная постоянная, то имеет место случай предельной точки (см. [19]), более общие результаты см. [20], [21].





\end{document}